\section{第二讲转场：从联想记忆到 Cognitive Map}

00:42:28 后课程更换讲者。原定嘉宾临时取消，Will Dorrell 在很短准备时间内用平板手写讲解，因此视觉形式从标准 slide deck 变成连续长页。内容却与前半紧密相连：Bricken 问“相似 query 怎样取回 memory”，Dorrell 问“经验中的关系结构怎样被记住，并在新内容上快速复用”。前者强调地址几何，后者强调结构与内容的 factorization。

\subsection{核心问题：序列经验中的结构能否被复用}

开场手写页把 hippocampus、entorhinal system 与 Transformers 放在一起。下一页用一段 sequence/graph 示意一个普遍问题：agent 观察到的对象和奖励会变化，但 transition structure、relative position 或 relational rule 可能保持不变。若模型每到新环境都从头记忆完整 episode，就无法 few-shot generalize；若能抽取 structure，就可把旧规则迁移到新内容。

\lecturefigure{slide-41-cognitive-map-title.jpg}{第二讲：Hippocampus、entorhinal system 与 Transformers}{官方录像手写教学状态 V041；00:43:54}

\lecturefigure{slide-42-one-problem-sequence-structure.jpg}{一个共同问题：从 sequence experience 抽取可复用结构}{官方录像手写教学状态 V042；00:45:36}

\teachervoice{课程主持人说明 Will 是临时接替，并介绍 Gatsby Unit 的 computational neuroscience 背景。Will 随即把自己的工作描述为：一个原本独立研究 hippocampal--entorhinal system 的模型，后来“看起来有点像 Transformer”。这句话强调的是 convergent computation，而非照着 Transformer 设计脑模型。}

\begin{importantbox}{Structural generalization 的最小定义}
给定不同 observation labels 但共享 transition/relational graph 的环境，系统应在少量新经验后复用已有结构进行 prediction、planning 或 inference。可迁移的是“关系怎样连接”，不是每个节点的感官身份。
\end{importantbox}

\subsection{本章小结}

第二讲把记忆问题从“相似度检索”推进到“结构复用”。认知地图不只要记住发生过什么，还要分离 what 与 where/how-related，使新的对象能够快速绑定到熟悉的关系图中。

